\documentclass[11pt,a4paper]{article}
\usepackage[margin=24mm]{geometry}
\usepackage{amsmath,amssymb,amsthm,booktabs,longtable,caption,hyperref}
\setlength{\emergencystretch}{2em}
\newtheorem{theorem}{Theorem}\newtheorem{proposition}[theorem]{Proposition}
\newcommand{\gen}[1]{\langle #1\rangle}\newcommand{\Aut}{\operatorname{Aut}}
\title{Finite coset geometries in Highly symmetric hypertopes}
\author{Human Compression project}\date{}
\begin{document}\maketitle
\begin{abstract}
We verify the three finite examples and all twenty-three table entries in
Fernandes--Leemans--Weiss, using their exact markings and presentations.
The common coset criteria are acquired from the earlier recovery of
\emph{Hexagonal extensions of toroidal hypermaps}; their application here
includes the actual flag and residue tests. All printed orders and
regularity assertions hold. One group name and both chiral flag-count
captions require correction. The completeness conjecture concerning the
chiral list is not used or proved.
\end{abstract}

\section{The groups and the certificate interface}
We use the accepted author manuscript of the source~\cite{source}, with
physical page numbers. Its Table~2, on page~20, specifies the presentations
\begin{equation}\label{hyp}
 P(p,a,b)=\gen{x,y,z\mid x^3,y^3,z^p,(x^{-1}z)^2,
 (y^{-1}z)^2,(x^{-1}y)^3,(xy^{-1}x)^a(xy)^b}.
\end{equation}
The involutory completion, when it exists, is part of the conclusion and
is not imposed in this presentation. Put $a_0=1,a_1=x,a_2=y,a_3=z$ and
\[
 P_J=\gen{a_i^{-1}a_j:i,j\in J},\qquad P_i=P_{\{0,1,2,3\}\setminus\{i\}}.
\]
The associated incidence system has right cosets of $P_i$ as its elements.
Two elements of different types are incident when their cosets intersect.

For Table~1 we use standard polytope rotations $u,v,w$. The literal
source relation (5.1) gives
\begin{equation}\label{poly}
 Q(p,a,b)=\gen{u,v,w\mid u^6,v^3,w^p,(uv)^2,(vw)^2,
 (uvw)^2,(vu^{-2})^a(v^{-1}u^2)^b}.
\end{equation}
Indeed $u=\rho_0\rho_1$, $v=\rho_1\rho_2$, $w=\rho_2\rho_3$, so
$vu^{-2}=\rho_1\rho_2(\rho_1\rho_0)^2$ and
$v^{-1}u^2=\rho_2\rho_1(\rho_0\rho_1)^2$.
The pair-difference marking for this group is
$(1,u,uv,uvw)$. In the program the three nonidentity entries are the free
generators; this is the invertible change $u=X,v=X^{-1}Y,w=Y^{-1}Z$.
It is essential to keep this change of marking when checking inversion.

We acquire the following complete finite certificate theorem from~\cite{finite}.
For four maximal rotation parabolics, set $K_J=\bigcap_{j\in J}P_j$.
The type-$i$ completions of the base flag of types $J$ are counted by
\[
 c(i,J)=\#\{g\in T:g\in P_iP_j\ (j\in J)\},
\]
where $T$ is a right transversal of $P_i\cap P_{j_0}$ in $P_{j_0}$ for
one $j_0\in J$. We check full pair-difference intersection, the equalities
$c(i,J)=[K_J:K_J\cap P_i]$ for $|J|=2$, and $c(i,J)=2$ for $|J|=3$.
In addition, the $P_i$ generate $P$, each $P_f$ is generated by its
intersections with the other parabolics, and every rank-two residue is
connected. The last test uses its actual incidence graph: two copies of
$K_J$, with $h$ joined to $k$ exactly when $hk^{-1}\in P_iP_j$, for the
two remaining types $i,j$.

These conditions imply a thin residually connected geometry with a free
action having exactly two chamber orbits. Thus they check more than
marked subgroup intersections. Simultaneous inversion of the three
marked generators is tested by evaluating \emph{all} defining relators
on their inverses. If the relators hold, the induced endomorphism squares
to the identity and is an automorphism. Otherwise a failed relator is a
literal obstruction. The acquired two-orbit theorem~\cite{chambers}
then determines the full type-preserving automorphism group. In the
regular case we also construct $H=P\rtimes\gen t$, put
$\rho_i=t a_i$, and verify all $256$ subgroup intersections and all
normalized flag-completion tests in $H$. The acquired regular-completion
proof identifies this geometry with the rotation coset geometry.

\section{All table entries}
\begin{theorem}\label{tables}
The presentations \eqref{hyp} and \eqref{poly}, at precisely the parameters
in Table~\ref{orders}, give the regular or chiral geometries recorded there.
The displayed numbers are the orders of their full type-preserving
automorphism groups. Each regular geometry has that many chambers;
each chiral geometry has twice that many chambers.
\end{theorem}
\begin{center}
\begin{tabular}{ccrrc}\toprule
$p$&$(a,b)$&Table 1 order&Table 2 order&status\\\midrule
3&(2,0)&240&120&regular\\
3&(3,0)&1296&648&regular\\
3&(4,0)&15360&7680&regular\\
3&(1,2)&336&168&chiral\\
3&(1,3)&2184&1092&chiral\\
3&(1,4)&8064&4032&chiral\\
3&(2,2)&2880&1440&regular\\
3&(2,3)&6840&3420&chiral\\
4&(1,1)&288&---&regular\\
4&(1,2)&2016&1008&chiral\\
4&(2,0)&768&384&regular\\
5&(2,0)&28800&14400&regular\\\bottomrule
\end{tabular}
\captionof{table}{The exact finite parameters.}\label{orders}
\end{center}
\begin{proof}
For every row, GAP computes an isomorphism from the finitely presented
group to a finite permutation group. This supplies both an upper bound
for the presentation and a faithful finite image; an arbitrary finite
quotient would not suffice. The complete seven-word relator lists are
retained. A separate literal-permutation checker verifies the relators,
generating orders and all pair-difference subgroup intersections without
calling the presentation routine.

The acquired finite criteria of Section~1 are then checked for every
row. There are $256$ pair-difference intersections, sixteen completion
counts, four parabolic-generation tests and six actual rank-two residue
graphs per row. The absent inversion is witnessed by the last relator
in each chiral row. Each regular row has the full involutory completion
and its independent intersection and flag checks. The proper-flag
counts prove flag completion and transitivity on all proper types;
the two completions of a panel lie in different free chamber orbits.
Consequently the full groups and chamber counts are those asserted.
For \eqref{poly} the string marking gives the corresponding abstract
polytope, by the standard string reconstruction theorem retained in
the general source recovery.

Local residue identifications are marked isomorphisms. In \eqref{hyp},
$\gen{x,y}$ is the exact torus group
\[
 \gen{x,y\mid x^3,y^3,(x^{-1}y)^3,(xy^{-1}x)^a(xy)^b}
\]
of order $3(a^2+ab+b^2)$. Its presentation maps onto the actual
parabolic, and independently computed equal orders prove injectivity.
The two designated spherical parabolics are checked in exactly the
same way against $\gen{s,t\mid s^3,t^p,(s^{-1}t)^2}$.
For \eqref{poly} the facet group has presentation
$\gen{u,v\mid u^6,v^3,(uv)^2,(vu^{-2})^a(v^{-1}u^2)^b}$,
of order $6(a^2+ab+b^2)$; the vertex group is the spherical
$\gen{v,w\mid v^3,w^p,(vw)^2}$.
The actual coset criteria already prove these groups describe residues;
the local names are not inferred solely from an unmarked order match.
\end{proof}

The statement about chambers corrects the captions of both original
tables. For example, Table~2 at $(p,a,b)=(3,1,2)$ has full automorphism
group of order $168$ and $336$ chambers. A ``flag'' in the table caption
must mean a chamber; even under that interpretation the printed $g$
is not the number of flags in a chiral row.

\section{The relation between the tables and the group names}
The group of each Table~2 row embeds with index two in the corresponding
Table~1 full group. This assertion is checked as a marked map, not by
comparing the two orders. In the rotations $(u,v,w)$ of \eqref{poly},
the map from \eqref{hyp} is
\begin{equation}\label{halving}
 x\longmapsto v^{-1}u^{-2},\qquad
 y\longmapsto v^{-1},\qquad z\longmapsto w.
\end{equation}
Every relator is satisfied, the image has half the order of $Q$, and
the computed homomorphism is bijective onto that subgroup. Interchanging
the first two images changes the orientation of the torus parameter;
we retain \eqref{halving} literally, including in the chiral rows.
In the regular completions the corresponding involution map is
\[
 (r_0,r_1,r_2,r_3)\longmapsto
 (\rho_2,\rho_1^{\rho_0},\rho_1,\rho_3).
\]
The omitted generator $\rho_0$ is an involution outside this index-two
subgroup. In the chiral cases an explicit involution outside the image
is likewise retained. These data verify the split doubling assertions,
including the unspecified actions in the source's semidirect-product
notation. No converse classification is inferred from this construction.

The symmetric, projective linear, direct-product and real Coxeter names
are checked by actual isomorphisms to their standard finite models.
The real Coxeter models use their complete reflection presentations.
The order-$648$ group is the monomial group $G(3,3,4)$:
\[
 \{(e_1,e_2,e_3,e_4)\in\mathbb F_3^4:\sum e_i=0\}\rtimes S_4,
\]
acting on four blocks of three points. An explicit isomorphism to this
model is retained. For the order-$7680$ group we use the named CHEVIE
model $G_{29}$~\cite{chevie}. Its four roots are the rows of
\[
 \frac12\begin{pmatrix}
 0&0&0&-2\\-i&i&1&1\\1-i&-1+i&0&0\\0&1-i&-1+i&0
 \end{pmatrix}.
\]
The reflections $v\mapsto v-2\langle v,r\rangle r/\langle r,r\rangle$
act on the orbit of these roots. Exact arithmetic in $\mathbb Q(i)$
gives $160$ roots and a permutation group of order $7680$. Since the
four roots span, that permutation action is faithful. The actual
isomorphism to the Table~2 group is checked. These explicit models remove
any reliance on the order alone for the source's abbreviations
$[1\,1\,2]^3$ and $[1\,1\,2]^4$. The historical abbreviations are retained
as labels; the displayed models specify the groups used here.

For $(p,a,b)=(3,1,4)$ the characteristic perfect residual has order
$336$ and is explicitly isomorphic to $\operatorname{SL}_2(7)$.
We retain a subgroup $A$ of order twelve, its isomorphism to $A_4$, and
the equalities $A\cap N=1$, $NA=G$. Thus the source's
$\operatorname{SL}_2(7)\rtimes A_4$ is proved with its actual conjugation
action. Its further extension in Table~1 is the split doubling above.

\begin{proposition}\label{correction}
The group in Table~2 with $(p,a,b)=(3,2,2)$ is
\[
 D=\{(s,t)\in S_5\times S_4:\operatorname{sgn}(s)=\operatorname{sgn}(t)\},
\]
not the printed $A_5\times S_4$.
\end{proposition}
\begin{proof}
On disjoint sets of five and four points, $D$ is generated by
\[
 (1\,2\,3),\ (3\,4\,5),\ (6\,7\,8),\ (7\,8\,9),\ (1\,2)(6\,7).
\]
These generate $A_5\times A_4$ and one element of odd sign on each
factor, hence the entire parity subgroup, of order $1440$.
The certificate gives a bijective homomorphism from the presented
Table~2 group to this model.

The stable last term of the derived series of $D$ is $A_5\times1$.
Indeed the first terms are $A_5\times A_4$, $A_5\times V_4$ and
$A_5\times1$. Its centralizer in $D$ is $1\times A_4$, of order twelve,
since $C_{S_5}(A_5)=1$. In contrast, the stable derived subgroup of
$A_5\times S_4$ is also $A_5\times1$, but its centralizer has order
$24$. This characteristic-subgroup invariant excludes an isomorphism.
The finite checker verifies both derived series and centralizer orders.
\end{proof}

\section{The three finite examples}
\begin{proposition}\label{examples}
Examples~4.4, 4.5 and 6.3 have all the group and incidence properties
asserted in the source.
\end{proposition}
\begin{proof}
For Example~4.4 put $u=ab$, $v=bc$, $t=b$, where $a,b,c$ are the three
printed involutions. Then $t$ inverts $u$ and $v$, and $(abc)^2=1$
is precisely $[u,v^{-1}]=1$. Thus the presentation is
$(C_3\times C_3)\rtimes C_2$, with simultaneous inversion, of order
$18$. Conversely this model satisfies every defining relation, so it
also supplies the lower bound. The three maximal parabolics have order
six, and any two intersect in their common involution. They therefore
have product of size $18$. Every pair of cosets of different types is
incident, giving $K_{3,3,3}$. All flags extend; there are $27$ chambers,
every panel has three completions, and every residue of rank at least
two is connected. Its full type-preserving group is $S_3^3$. In
particular the designated order-$18$ group is not flag transitive.

For Example~4.5 use the literal generators
\[
 (12)(34),\quad(26)(35),\quad(14)(23),\quad(14)(35).
\]
They generate $A_6$ and satisfy all $256$ marked subgroup intersections.
Their maximal parabolics have orders $(24,60,60,24)$ and the four
type sets have sizes $(15,6,6,15)$. The complete coset calculation gives
$381$ incidence edges, $1080$ three-type flags and $720$ chambers.
All shorter flags extend. Among the $1080$ panels, $720$ have two
completions and $360$ have four. A residue of types $1,2$ through a
base flag of types $0,3$ is $K_{4,4}$. A base flag of types $1,2$ has
a residue consisting of two disjoint $10$-cycles. There are precisely
$36$ disconnected residues of rank at least two. Hence thinness,
residual connectivity and designated flag transitivity all fail.
The last assertion also follows from $|A_6|=360<720$.

For Example~6.3 retain exactly the four permutations printed on page~14.
Their group has order $72$, with marked subgroup orders, in binary
subset order,
\[
 1,2,2,6,2,4,8,24,2,6,6,18,4,12,24,72.
\]
All $256$ intersections are correct. The rotation subgroup is the
whole group. For $I=\{0,1,2\}$ and $J=\{1,2,3\}$,
\[
 |P_I|=|P_J|=24,\qquad |P_I\cap P_J|=8,
 \qquad |P_{I\cap J}|=4.
\]
The source's second involution $\rho_1$ belongs to the first
intersection and not to $P_{I\cap J}$, giving a literal witness.
These computations are independently repeated by closure of the
permutation arrays and set intersections, without subgroup or coset
routines from GAP. For Examples~4.4 and 4.5 this independent check also
constructs every coset, every flag and every residue incidence graph.
\end{proof}

\section{Scope, parameters and attribution}
A minor domain qualification is needed in the source's rank-three
toroidal discussion. Its unrestricted wording excludes $(1,1)$, but
one must also exclude the degenerate unit periods. For example,
\[
 \gen{x,y\mid x^3,y^3,(x^{-1}y)^3,xy^{-1}x}
 =\gen{x\mid x^3}=C_3,
\]
since $y=x^2$. Its three pair-difference parabolics are all the whole
group; its standard coset system has one element of each type and is
not thin. The usual admissible parameter domain is therefore necessary
when invoking the prior toroidal classification. We make no claim here
that a finite sampling proves that classification. Every period in
Table~\ref{orders} has norm at least three, and every Table~2 period has
norm at least four; all are checked separately.

The source attributes the regular classification of locally toroidal
polytopes to McMullen--Schulte, the rank-three chirality criterion to
Breda D'Azevedo--Nedela, and the toroidal hypermap families to
Corn--Singerman. Those classification statements remain explicit prior
inputs wherever the general source recovery uses them. Our two finite
tables are now verified positive examples independently of their
completeness. Problem~9.2 and the unnumbered completeness conjecture
remain separate questions. The old negative answers to Problems~9.3
and 9.4, including their restricted Lean scope, are acquired in the
general recovery and receive no new credit here.

The source presentations, examples and positive conclusions belong to
the source authors and their cited predecessors. The finite coset
certificate theorem and its original implementation are reused from the
accepted L003 recovery. The present contribution consists of the literal
L052 applications, independent permutation and incidence checks, exact
named-group maps, and the explicit corrections above. Named finite-group
constructors, exact GAP presentation routines, and the named CHEVIE root
model are declared computational and identification premises. They are
not hidden inside a claim of Lean verification or new classification.

\begin{thebibliography}{9}\small
\bibitem{source} M. E. Fernandes, D. Leemans and A. Ivi\'c Weiss,
\emph{Highly symmetric hypertopes}, Aequationes Math. 90 (2016),
1045--1067; accepted author manuscript, fixed 22-page edition.
\bibitem{finite} \emph{Exact rotation presentations and coset certificates}
and \emph{Literal local periods and actual marked quotients}, accepted
common proofs in the complete L003 recovery, Sections 00--01 of its
shared supplement; exact editions and acceptance records acquired.
\bibitem{chambers} \emph{Chamber propagation and two regular orbits},
accepted Topic Chamber common manuscript, version 2.
\bibitem{chevie} G. Malle and J. Michel, CHEVIE data for $G_{29}$,
\texttt{src/tbl/cmplxg29.jl}; four literal root vectors, retrieved from
the authors' software repository. See J. Michel, J. Algebra 435 (2015),
308--336, for the package.
\end{thebibliography}
\end{document}
